\section*{Capítulo \ref{ch:g6:producing-our-food} --- Produzir o nosso alimento}

\begin{solution}{exo:g6:producing-our-food:1}
O \omterm{def:g6:producing-our-food:farming}{cultivo} maneja produtoras --- \omterm{def:g1:plants-around-us:plant}{plantas} cultivadas pela \omterm{def:g6:plants-are-producers:matter}{matéria
orgânica} delas. A \omterm{def:g6:producing-our-food:farming}{criação de animais} maneja \omterm{def:g5:food-webs:roles}{consumidores} --- \omterm{def:g1:animals-around-us:animal}{animais}
criados pela produção deles: \omterm{def:g2:animals-grow-up:born}{leite}, carne, \omterm{def:g2:animals-grow-up:born}{ovos}, lã.
\end{solution}

\begin{solution}{exo:g6:producing-our-food:2}
Grão: a \omterm{def:g2:from-seed-to-plant:seed}{semente} do pé de trigo (a reserva embalada dele). Batata: o
\omterm{ex:g4:plants-without-seeds:bulbtuber}{tubérculo} do pé de batata. Maçã: o \omterm{prop:g3:flowers-fruits-seeds:transform}{fruto} da macieira. \omterm{def:g2:animals-grow-up:born}{Leite}: o produto
do cuidado materno da vaca. \omterm{def:g2:animals-grow-up:born}{Ovo}: o berçário embalado da galinha.
\end{solution}

\begin{solution}{exo:g6:producing-our-food:3}
A transformação de uma matéria-prima pela alimentação de \omterm{def:g6:producing-our-food:fermentation}{micróbios}. O
pão (matéria-prima: a massa de farinha; trabalhador: o \omterm{ex:g6:producing-our-food:bread}{fermento}) e o
iogurte (matéria-prima: o \omterm{def:g2:animals-grow-up:born}{leite}; trabalhadores: os \omterm{def:g6:producing-our-food:fermentation}{micróbios} que
azedam); o queijo e o vinagre também servem.
\end{solution}

\begin{solution}{exo:g6:producing-our-food:4}
O \omterm{ex:g6:producing-our-food:bread}{fermento} --- um \omterm{def:g6:decomposers-and-soil:decomposers}{fungo} de uma \omterm{def:g5:first-look-at-cells:cell}{célula} só --- se alimenta dos açúcares
da farinha e libera um gás; a massa elástica prende esse gás em bolhas
e incha. O calor do forno encerra o trabalho do \omterm{ex:g6:producing-our-food:bread}{fermento} e fixa o
miolo crescido, com buraquinhos e tudo.
\end{solution}

\begin{solution}{exo:g6:producing-our-food:5}
O ácido liberado pelos \omterm{def:g6:producing-our-food:fermentation}{micróbios} que se alimentam do açúcar do \omterm{def:g2:animals-grow-up:born}{leite}:
o ácido engrossa e coalha o \omterm{def:g2:animals-grow-up:born}{leite}, e o azedinho fresco é o próprio
ácido.
\end{solution}

\begin{solution}{exo:g6:producing-our-food:6}
Salgar, resfriar (geladeira), secar --- e também cozinhar: cada um
nega aos \omterm{def:g6:producing-our-food:fermentation}{micróbios} indesejados uma condição de trabalho (\omterm{def:g6:exploring-our-environment:conditions}{umidade},
calor), exatamente como o frio deixa lenta uma pilha de composto.
\end{solution}

\begin{solution}{exo:g6:producing-our-food:7}
Água: quem cultiva conta com a chuva (e às vezes a leva até lá). \omterm{def:g6:exploring-our-environment:conditions}{Luz}:
plantio em campo aberto, com espaçamento para cada \omterm{def:g1:plants-around-us:plant}{planta} receber a
parte dela. Calor: plantio no tempo da estação quente. Ar: terreno
aberto, \omterm{def:g6:decomposers-and-soil:soil}{solo} afofado. \omterm{prop:g4:what-plants-need:needs}{Sais minerais}: esterco e reabastecimento depois
da colheita.
\end{solution}

\begin{solution}{exo:g6:producing-our-food:8}
(1) Organismo e etapa: o \omterm{def:g2:animals-grow-up:born}{leite} da vaca --- produto de um \omterm{def:g5:food-webs:roles}{consumidor}.
(2) Lugar no ciclo: produto de \omterm{def:g5:food-webs:roles}{consumidor}, acabado por \omterm{def:g6:producing-our-food:fermentation}{micróbios}.
(3) Matéria-prima: o \omterm{def:g2:animals-grow-up:born}{leite}; os trabalhadores o azedam e o coalham em
coalhada e depois curam o queijo prensado por meses --- buracos, casca
e cheiro são o registro deles. (4) Trilha: o \omterm{def:g2:animals-grow-up:born}{leite} veio da vaca, a
vaca veio do capim, o capim veio da \omterm{def:g6:exploring-our-environment:conditions}{luz}.
\end{solution}

\begin{solution}{exo:g6:producing-our-food:9}
Os \omterm{def:g6:producing-our-food:fermentation}{micróbios} são trabalhadores vivos com \omterm{def:g6:exploring-our-environment:conditions}{condições}: o frio quase para
a alimentação deles. Na geladeira isso protege o alimento (os
estragadores ficam ociosos); na pilha de inverno isso atrasa o
composto (os \omterm{def:g6:decomposers-and-soil:decomposers}{decompositores} ficam ociosos). O princípio: a
\omterm{prop:g6:decomposers-and-soil:decomposition}{decomposição} e o estrago acontecem no ritmo que as \omterm{def:g6:exploring-our-environment:conditions}{condições}
permitem.
\end{solution}

\begin{solution}{exo:g6:producing-our-food:10}
``\omterm{def:g1:healthy-habits:germs}{Micróbio}'' é o nome de uma multidão enorme de \omterm{def:g6:exploring-our-environment:components}{seres vivos} minúsculos
e diferentes entre si. Alguns, no lugar errado, causam doença --- são
esses o alvo da lavagem das mãos; outros, escolhidos e manejados em
\omterm{def:g6:exploring-our-environment:conditions}{condições} controladas, fazem o pão crescer e o \omterm{def:g2:animals-grow-up:born}{leite} coalhar. Inimigo
ou funcionário não é a natureza do \omterm{def:g1:healthy-habits:germs}{micróbio}, e sim a combinação de
\omterm{def:g5:biodiversity:species}{espécie}, lugar e manejo.
\end{solution}

\begin{solution}{exo:g6:producing-our-food:11}
O presunto passou por um porco: o peso dele custou muitas vezes esse
peso em ração, mais a fatia queimada do porco, mais terra e água --- o
pedágio do andar dos \omterm{def:g5:food-webs:roles}{consumidores}. O trigo do pão veio a um andar da
\omterm{def:g6:exploring-our-environment:conditions}{luz}. Desperdiçar o sanduíche desperdiça os dois balanços; o do
presunto é o mais comprido.
\end{solution}

\begin{solution}{exo:g6:producing-our-food:12}
Eles manejavam as \emph{\omterm{def:g6:exploring-our-environment:conditions}{condições}} dos trabalhadores: o calor do canto
da massa, o \omterm{ex:g6:producing-our-food:bread}{fermento} natural guardado e alimentado, o sal e os tempos
ajustados por tentativa ao longo de gerações. As receitas funcionavam
porque manter as \omterm{def:g6:exploring-our-environment:conditions}{condições} certas mantém prosperando os \omterm{def:g6:exploring-our-environment:components}{seres vivos}
certos --- saiba-se ou não que os \omterm{def:g6:exploring-our-environment:components}{seres vivos} estão ali. Eles criavam
\omterm{def:g6:producing-our-food:fermentation}{micróbios} às cegas, e criavam bem.
\end{solution}

\begin{solution}{pb:g6:producing-our-food:1}
\textbf{1.} Cada grão é uma \omterm{def:g2:from-seed-to-plant:seed}{semente}: um pé de trigo minúsculo embalado
com uma \omterm{def:g3:animals-through-seasons:active}{reserva de alimento}, feita pela planta-mãe para os primeiros
dias daquela plantinha --- não para nós, embora a gente a desvie.

\textbf{2.} Nas próprias \omterm{def:g1:plants-around-us:parts}{folhas} verdes dos pés de trigo, a partir de
água e \omterm{prop:g4:what-plants-need:needs}{sais minerais} do \omterm{def:g6:decomposers-and-soil:soil}{solo} e do gás carbônico do ar, com a \omterm{def:g6:exploring-our-environment:conditions}{luz} como
\omterm{prop:g3:balanced-diet:jobs}{energia}.

\textbf{3.} Plantado, um grão percorre o \omterm{def:g2:born-grow-die:cycle}{ciclo de vida} dele: da
\omterm{def:g2:from-seed-to-plant:germination}{germinação} à \omterm{def:g1:plants-around-us:plant}{planta} e a grãos novos. Vendido e moído, o \omterm{def:g2:born-grow-die:cycle}{ciclo de vida}
dele acaba e a \omterm{def:g3:animals-through-seasons:active}{reserva de alimento} dele vira nossa --- um lanche
embalado para uma plantinha, comido por uma pessoa.

\textbf{4.} Devolver o que se toma emprestado: o esterco, ao
apodrecer, reabastece os \omterm{prop:g4:what-plants-need:needs}{sais minerais} que a colheita levou embora ---
servindo à necessidade de \omterm{prop:g4:what-plants-need:needs}{sais minerais}.

\textbf{5.} Porque o \omterm{ex:g6:producing-our-food:bread}{fermento} é um trabalhador vivo: no calor, ele se
alimenta e trabalha rápido; no frio, fica ocioso --- a regra da equipe
do composto, em escala de padaria.

\textbf{6.} Ele se alimentou dos açúcares da farinha e liberou gás; a
massa elástica prendeu o gás em bolhas, e as bolhas dobraram a massa.

\textbf{7.} A janela fria deixou o \omterm{ex:g6:producing-our-food:bread}{fermento} ocioso --- quase sem se
alimentar, quase sem gás. As duas equipes fizeram sem querer um teste
justo do calor: a mesma massa, o mesmo dia, uma condição diferente,
dois crescimentos para comparar.

\textbf{8.} O calor matou o \omterm{ex:g6:producing-our-food:bread}{fermento} e encerrou o turno; o miolo
firme guarda os buraquinhos das bolhas --- a assinatura da força de
trabalho, assada no lugar.

\textbf{9.} A fatia vem do pão, o pão da massa, a massa da farinha, a
farinha do grão, o grão das \omterm{def:g1:plants-around-us:parts}{folhas} do pé de trigo --- fabricado ali a
partir de ingredientes minerais com a \omterm{def:g6:exploring-our-environment:conditions}{luz}: a trilha termina no sol da
\omterm{def:g6:producing-our-food:farming}{lavoura} de julho.

\textbf{10.} Fatia: um andar (da produtora à mesa, com um acabamento
de \omterm{def:g6:producing-our-food:fermentation}{micróbios}). Manteiga: dois andares (capim, vaca, creme). Presunto:
dois andares, com o maior pedágio --- o porco queimou e eliminou quase
toda a matéria da ração dele antes de qualquer parte virar presunto. A
fatia é a mais barata em matéria; o presunto, a mais cara.

\textbf{11.} Os \omterm{def:g6:decomposers-and-soil:decomposers}{decompositores} da composteira --- minhocas,
tatuzinhos, \omterm{def:g6:decomposers-and-soil:decomposers}{fungos}, equipes invisíveis --- desmontam a \omterm{def:g6:plants-are-producers:matter}{matéria
orgânica} da massa em \omterm{def:g6:decomposers-and-soil:soil}{húmus} e \omterm{prop:g4:what-plants-need:needs}{sais minerais}: mais uma volta, do
alimento para a fertilidade do \omterm{def:g6:decomposers-and-soil:soil}{solo} do ano que vem.

\textbf{12.} Por exemplo: ``O nosso pão foi construído por uma
produtora a partir de \omterm{def:g6:plants-are-producers:matter}{matéria mineral} e de \omterm{def:g6:exploring-our-environment:conditions}{luz}, foi levantado por uma
força de trabalho de \omterm{def:g6:producing-our-food:fermentation}{micróbios} e não precisou de andar nenhum de
\omterm{def:g5:food-webs:roles}{consumidor} --- embora os sanduíches de presunto do lado precisassem de
um.''
\end{solution}
